\section{Conclusiones}

\begin{itemize}
\item Los resultados obtenidos permitieron comprobar experimentalmente que una red resistiva lineal de corriente continua puede representarse mediante su circuito equivalente de Thévenin, conservando el comportamiento eléctrico observado en los terminales de interés. Esto permitió comprender de manera práctica la utilidad de este teorema como herramienta para simplificar el análisis de circuitos.

\item La comparación entre los valores teóricos y experimentales del voltaje y la resistencia de Thévenin mostró una correspondencia elevada, con errores relativos inferiores al $1\%$ en la mayoría de las mediciones. La mayor diferencia registrada fue del $2\%$ en el primer ensayo del segundo circuito, lo que evidencia que las discrepancias experimentales fueron pequeñas en relación con los valores esperados.

\item La corriente obtenida al conectar la resistencia de carga presentó errores relativos entre el $1\%$ y el $3\%$ respecto a los valores teóricos. Estos resultados muestran que, aunque las mediciones de corriente presentan una mayor sensibilidad a las condiciones experimentales y a la resolución del instrumento, el comportamiento del circuito equivalente continúa siendo consistente con el circuito original.

\item Las diferencias entre los resultados teóricos y experimentales evidencian la influencia de las condiciones reales de medición, tales como la resolución del multímetro, la incertidumbre asociada a las lecturas y las tolerancias de los componentes. Por tanto, los resultados muestran la importancia de considerar las incertidumbres experimentales al comparar un modelo teórico con un sistema físico real.

\item En conjunto, los resultados obtenidos presentan un grado de concordancia suficientemente alto para validar experimentalmente el Teorema de Thévenin bajo las condiciones estudiadas. El experimento permitió comprender que la utilidad del teorema no se limita a una simplificación matemática, sino que constituye una herramienta práctica para analizar el comportamiento de una carga conectada a una red eléctrica lineal.


\end{itemize}


